\section{Related Work}
\paragraph{LLM Reasoning.} As an emerging capability of model scale, LLMs can generate intermediate CoTs to solve complex reasoning tasks \citep{wei2022chain,kojima2022large} and scale test-time performance by allocating more thinking tokens \citep{snell2024scaling,brown2024large}. Early efforts enhanced LLM reasoning via supervised fine-tuning on human-annotated data \citep{cobbe2021training,yue2023mammoth,yu2023metamath} or linearized search traces \citep{lehnert2024beyond,gandhi2024stream}. 
However, due to the distribution shift between LLM responses and curated data, LLM-generated data has proven effective through rejection sampling \citep{dong2023raft,yuan2023scaling} by filtering out low-quality rationales \citep{zelikman2022star,zelikman2024quiet} or with EM iterations \citep{singh2023beyond,zhong2025brite}. Recently, RL has gained increasing interest for improving reasoning \citep{aksitov2023rest,havrilla2024teaching,wang2024offline,shao2024deepseekmath}. Process rewards \citep{uesato2022solving,lightman2023let} with Monte Carlo unrolls \citep{kazemnejad2024vineppo,wang2023math,wang2024multi,luo2024improve} offer finer-grained feedback but are computationally expensive. Outcome-reward RL \citep{guo2025deepseek} demonstrates emergent deliberative reasoning abilities such as self-reflection. Yet, limited work has investigated the underlying mechanisms of such behaviors. In fact, recent findings suggest that reflections do not consistently emerge from RL training and exhibit weak correlation with performance \citep{liu2025understanding}.
Similar to our work, \cite{xiang2025towards,qu2025optimizing} also study reflective exploration of LLMs, from a meta-RL \citep{duan2016rl} perspective: \cite{xiang2025towards} justifies deliberative reasoning as providing extra CoT state contexts, and \cite{qu2025optimizing} uses progress reward \citep{setlur2024rewarding} to reduce regret in outcome-reward RL. Our method differs from \cite{xiang2025towards} in that we ground reflective reasoning in environment rewards, rather than relying solely on the internal CoT states generated by the model itself. Compared to \cite{qu2025optimizing}, which rewards strategies that make progress towards the correct answer, BARL additionally encourages exploring plausible strategies under the Bayesian framework. We defer a more detailed discussion to the end of Section \ref{eq_nece}. We experimentally compare with a variant of \cite{qu2025optimizing} that estimates progress using answer probability differences. 
Prior work has also brought Bayesian exploration principles to LLMs: \cite{arumugam2025toward} implements \emph{Posterior Sampling for Reinforcement Learning} (PSRL), a statistically-efficient algorithm, using LLMs as modular components, while \cite{dwaracherla2024efficient} shows that active query selection with epistemic neural networks and double Thompson sampling reduces the amount of human feedback required to train reward models for RLHF. 
In contrast, we optimize a Bayesian objective directly during RL fine-tuning of the LLM policy, rather than at the level of external algorithmic scaffolding or data collection, to provide a step-level guidance on when and how LLMs should self-reflect.
Besides, our method differs from \cite{wang2023hypothesis,qiu2023phenomenal} that manually designs hypothesis proposal–selection pipelines.

\vspace{-0.15cm}
\paragraph{Reinforcement Learning.} Conventional RL explores only during training. Exceptions include works that explicitly optimize maximum entropy objectives \citep{haarnoja2018soft} to learn stochastic policies, primarily to accelerate \textit{training} convergence in settings where evaluation remains in-distribution, such as robotic control. Bayes-Adaptive RL \citep{bellman1959adaptive,duff2002optimal,guez2012efficient,ghavamzadeh2015bayesian,lidayan2025bamdp,liu2023reason} has been studied to pursue the optimal exploration-exploitation trade-off in uncertain environments.
When the true MDP identity is latent and must be inferred from interaction (states, actions, rewards), Bayesian RL connects naturally to Partially Observable MDPs \citep{duff2001monte,ghosh2021generalization}. Exact solutions of the Bayesian RL objective are often intractable, promoting the development of approximate methods \citep{guez2012efficient,arumugam2022planning,chen2024bayes,arumugam2024satisficing}. In our work, we adopt a policy gradient that operates over candidate answers, which differs from \cite{ghosh2022offline} that leverages value ensembles in offline RL and \cite{qiu2025bayesian} that applies SFT on an oracle Bayesian model's outputs. 
